% Ad Atticum 15.20 (ad-atticum-15-20) --- English translation
% Translated by Claude (Anthropic) from the Latin (Perseus).
% Style guide: STYLE.md.

\ciceroLetterOpener{CICERO ATTICO salutem}

\ciceroSection{1} I have thanked Vettienus; nothing could have been more courteous. Let Dolabella's instructions be what they will, let me have something --- if only what I can report to Nicias. For who, as you write, \dag is hesitating now\dag --- who that is prudent now doubts that my leaving is an act of despair, not of embassy? As to what you say, that men are now beginning to talk extremes about the state, and good men at that: from the day on which I heard that tyrant called, in the public assembly, ``the most illustrious of men,'' I began to lose confidence. After I saw, with you at Lanuvium, that our own side had only as much hope of staying alive as they had received from Antony, I despaired. And so, my dear Atticus --- take this resolutely, as I write it --- judging foul that mode of death which \dag I am bound to fall into\dag, and as it were announced to us by Antony, I have determined to get out of this lobster-pot, not in flight, but in the hope of a better death. The whole of the blame is Brutus's.

\ciceroSection{3} You write that Pompeius has been received at Carteia: there is therefore an army now in the field against this man. Which camp, then? For Antony does away with any middle ground. The one is weak, the other criminal. Let us press on, then. But help me with your counsel: Brundisium or Puteoli? Brutus, indeed, went off suddenly --- but wisely. I am suffering somewhat \textit{[Greek: pasch\=o ti]}. For when shall I see him again? Yet human things must be borne. You yourself cannot see him. Curse the dead man who ever made Buthrotum the issue! But what is done is past; let us see what is to be done.

\ciceroSection{4} Eros's accounts, though I have not yet seen him in person, I have nevertheless got fairly fully from his letter and from what Tiro has learned. You write that a loan will have to be raised for five months, that is to the Kalends of November, of 200,000 sesterces, and that monies due from Quintus fall in by that date. Since, then, Tiro says you do not approve of my coming to Rome for the purpose, I should like you, if it gives you no offence, to see where the money is to come from, and to charge it to me. This, I see, is what is needed for the present. The rest I shall make out more carefully from the man himself, among other things about the rents from my wife's dower-estates. If those can be faithfully managed for young Cicero, although I wish to provide more amply, even so he will lack for almost nothing. I see that I myself also need travelling-money; but for him, as the income falls due from the estates, so it will be paid, whereas for me I need the full sum at once. For my part, although the man who fears shadows seems to me to be aiming at bloodshed, still I shall not leave until the financial business is sorted out. Whether it is sorted out or not I shall find out from you. I thought this had better be written in my own hand, and so I did. About Fadius, as you write, by all means to no one else. I should be glad if you wrote back today.
